\section{Finite presentations and observation contracts}

\subsection{Presentations}

The framework starts before proof syntax.  A mathematical object must first be
replaced by a finite-use presentation.

\begin{definition}[Finitary presentation]
A finitary presentation \(X\) consists of primitive-recursive data
\[
 (\mathsf{Code}_X,\mathsf{Query}_X,\mathsf{Obs}_X,
   \mathsf{valid}_X,\mathsf{restrict}_X)
\]
with the following intended meanings.
\begin{enumerate}
\item \(x\in\mathsf{Code}_X\subseteq\N\) is a code or name for an input.
\item \(q\in\mathsf{Query}_X\) is a finite request: precision, radius, degree,
finite subdiagram, support bound, or resource level.
\item \(o\in\mathsf{Obs}_X(x,q)\) is a finite observation code.
\item \(\mathsf{valid}_X(x,q,o)\in\{0,1\}\) is a decidable validity predicate.
\item If \(q\preceq q'\), the restriction operation maps a valid observation at
\(q'\) to one at \(q\).
\end{enumerate}
All domains may be represented by natural-number bounds and characteristic
terms; the typed notation is explanatory.
\end{definition}

The coherence law is
\[
  q\preceq q'\preceq q'',\quad
  \mathsf{valid}(x,q'',o'')
  \quad\Longrightarrow\quad
  \mathsf{restrict}_{q'',q}(o'')
  =\mathsf{restrict}_{q',q}
      (\mathsf{restrict}_{q'',q'}(o'')).
\]
This law is an equation between PR terms or a bounded formula, hence suitable
for the object theory.  It prevents the presentation from offering unrelated
answers at different precisions.

A presentation need not have a unique observation.  This is essential.  An
algebraic number may be represented by an isolating disk; a root at a collision
may have several compatible disks; a compact space may have many finite nets;
a quotient object may have many representatives.  The validity relation is
therefore primary, and a deterministic name is a special case.

\subsection{Deterministic and relational names}

\begin{definition}[Deterministic name]
A deterministic name is a primitive-recursive function
\(a_X(x,q)\) with a proof of
\(\mathsf{valid}_X(x,q,a_X(x,q))\) and restriction coherence.
\end{definition}

\begin{definition}[Finite relational name]
A finite relational name supplies a primitive-recursive bound \(B_X(x,q)\) and
a validity predicate such that
\[
        \exists o\le B_X(x,q)\;\mathsf{valid}_X(x,q,o).
\]
The cloud at \(q\) is the finite set of all valid \(o\le B_X(x,q)\).
\end{definition}

The existential statement may itself be produced constructively, or it may be
a hypothesis supplied by a deep theorem.  Once the bound and validity relation
are available, complete enumeration, least search, and all universal checks
over the cloud are primitive recursive.

\subsection{Observation contracts}

An operation is specified by what it must observe, not by an unbounded semantic
function.

\begin{definition}[Deterministic observation contract]
A deterministic contract \(F:X\to Y\) consists of:
\begin{enumerate}
\item a backward demand transformer
\(d_F(x,r)\in\mathsf{Query}_X\);
\item a primitive-recursive transformer
\(f_F(x,r,o)\in\mathsf{Obs}_Y\);
\item a local proof that
\[
 \mathsf{valid}_X(x,d_F(x,r),o)
 \longrightarrow
 \mathsf{valid}_Y(Fx,r,f_F(x,r,o)).
\]
\end{enumerate}
The displayed implication is ultimately encoded as a native clause or ordinary
iPRA formula.
\end{definition}

\begin{definition}[Relational observation contract]
A relational contract replaces \(f_F\) by a finite output bound and relation
\(R_F(x,r,o,v)\).  Its local proof states that every valid demanded input
observation has at least one valid related output and that every related output
satisfies the advertised \(Y\)-contract.
\end{definition}

A contract is not a semantic assertion that a host function is continuous.  It
is a finite theorem: a particular input observation at a computed demand
suffices to verify a particular output observation.

\subsection{Composition of deterministic contracts}

Let \(F:X\to Y\) and \(G:Y\to Z\) be deterministic contracts.  For a requested
query \(s\) on \(Z\), first compute the \(Y\)-query
\(r=d_G(Fx,s)\), then the \(X\)-query
\[
     d_{G\circ F}(x,s)=d_F(x,r).
\]
Given an observation \(o\) at this demand, form
\[
   u=f_F(x,r,o),\qquad v=f_G(Fx,s,u).
\]
The local proof for \(F\) gives validity of \(u\), and the local proof for
\(G\) gives validity of \(v\).  In the proof compiler the corresponding
operation is cut.

\begin{theorem}[Deterministic contract composition]
Deterministic observation contracts are closed under identity and composition.
Demand transformers and realizers remain primitive recursive, and local
object-language proofs compose by cut.
\end{theorem}

\begin{proof}
The identity demand is the requested query and the identity realizer is
restriction or literal identity.  For composition, the formulas above are
primitive-recursive compositions.  The proof of the composite contract is the
first local derivation followed by the second, discharging the intermediate
validity formula.  The formal version is Theorem~\ref{thm:list-cut} below.
Associativity holds extensionally because primitive-recursive composition and
proof cut are associative up to derivability; no equality of proof terms is
needed.
\end{proof}

\subsection{Adaptive finite use}

Many algorithms do not know all input requests in advance.  A bisection routine
chooses its next interval from a sign test; a pivoting algorithm asks for a gap
after inspecting a candidate pivot; a continuation method changes precision
when an error estimate is large.  Such an algorithm is still finitely
presentable when it is clocked.

A state \(s_t\) contains both mathematical data and a finite request log.  At
step \(t\), a primitive-recursive rule computes the next query from \(s_t\),
reads the corresponding input observation, and produces \(s_{t+1}\).  For a
fixed fuel \(T\), the total demand is the maximum of the finitely many queries
recorded in the trace.  Rather than pretending that the adaptive algorithm is a
single nonadaptive contract, the framework compiles its trace and proves the
invariant that every query was answered correctly.

\subsection{Products, sums, subobjects, and quotients}

Finite presentations are closed under the ordinary constructions needed by
mathematics.

\paragraph{Products.}
A query for \(X\times Y\) is a pair of queries, an observation is a pair of
observations, and validity is conjunction.  Pairing functions code the pair by
a natural number.  Projections are exact PR operations.

\paragraph{Finite sums.}
An observation has a tag and a payload.  Validity checks the tag and invokes
the relevant component predicate.  Exact case analysis is permitted because
the tag is a decidable natural number.

\paragraph{Subobjects.}
A subpresentation is a presentation plus a decidable or explicitly certified
property.  The property is carried as part of validity.  For an open property,
one must decide whether the theorem uses it as a hypothesis or supplies a
finite certificate.

\paragraph{Quotients.}
A quotient should generally not choose canonical representatives.  It is
presented by representatives, a bounded equivalence relation on the demanded
finite data, and observers invariant under that relation.  A normal form may be
added when a primitive-recursive normalization theorem is available, but it is
not part of the general quotient interface.

\subsection{Adequacy without reflection}

Each presentation has an external adequacy theorem connecting its finite codes
to the intended mathematical object.  For represented reals, adequacy says
that a coherent rapidly Cauchy rational name determines a real.  For finite
simplicial complexes, the code decodes to an abstract complex.  For matrices,
the code decodes to a fixed finite table.

Adequacy is not used as a truth-to-provability rule.  The object theorem is
stated directly about codes and validity predicates.  Adequacy explains why
that theorem is the desired mathematical statement.  This separation permits
classical mathematics to be used in a metatheoretic interpretation without
silently importing classical axioms into the derivation.

\subsection{Producers, verifiers, and finalizers}

A recurring source of overclaim is failure to separate three roles.

\begin{definition}[Producer--verifier--finalizer factorization]
A theorem pipeline is factored as follows.
\begin{enumerate}
\item A \emph{producer} supplies candidate finite data and quantitative bounds.
It may be a primitive-recursive algorithm, an explicit theorem hypothesis, or a
separately established deep result.
\item A \emph{verifier} is a primitive-recursive bounded checker for the local
contract.
\item A \emph{finalizer} turns verified finite data into the stated output and
carries the object-language derivation of correctness.
\end{enumerate}
\end{definition}

The general framework proves closure of verifiers and finalizers.  It proves a
producer only when an actual mathematical construction and its quantitative
analysis are supplied.  In particular, saying ``let the certificate contain a
height bound'' does not prove a height bound.
